\section{PI Controller Pacing}
\label{sec:exp4b_results}

Experiment~4b replaces the multiplicative dual pacing of Experiment~4a with a PI controller. The $2^6$ factorial crosses auction format, aggressiveness ($a_{\text{low}}$ vs.\ $a_{\text{high}}$), market thickness ($N = 2$ vs.\ $N = 4$), budget multiplier ($m = 0.25$ vs.\ $m = 1.0$), reserve price ($r = 0.0$ vs.\ $r = 0.3$), and value dispersion ($\sigma = 0.1$ vs.\ $\sigma = 0.5$), each replicated across 8 independent seeds (512 runs total).

\begin{figure}[H]
  \centering
  \includegraphics[width=\textwidth]{figures/e4b_trace}
  \caption{Representative learning trajectory for a single trial of Experiment~4b (first-price auction, default aggressiveness, 2~bidders, 100 episodes $\times$ 1{,}000 rounds). Mean bid-to-value ratio per episode with the competitive benchmark ($b/v = 1.0$).}
  \label{fig:e4b_trace}
\end{figure}

\subsection{Efficiency and Welfare}

The PI controller produces welfare outcomes that can be compared directly with the dual pacing results of Experiment~4a. Platform revenue and liquid welfare are computed identically, and the effective Price of Anarchy uses the same LP-based offline optimum as denominator. The budget multiplier again dominates revenue (Table~\ref{tab:exp4b_ranked_rev}), with the effect amounting to \ExpFourBRevTopPctEffect\% of the grand mean, confirming the pattern from Experiment~4a.

\input{tables/exp4b_ranked_rev}

\begin{figure}[H]
  \centering
  \includegraphics[width=0.7\textwidth]{figures/e4b_main_rev}
  \caption{Experiment~4b: Main effects plot for platform revenue.}
  \label{fig:e4b_main_rev}
\end{figure}

\input{tables/exp4b_ranked_poa}

\begin{figure}[H]
  \centering
  \includegraphics[width=0.7\textwidth]{figures/e4b_main_poa}
  \caption{Experiment~4b: Main effects plot for effective Price of Anarchy.}
  \label{fig:e4b_main_poa}
\end{figure}

For the effective Price of Anarchy, auction format is the dominant factor ($|t| = 10.7$), with first-price auctions producing higher PoA values. The number of bidders $\times$ budget multiplier interaction ranks second ($|t| = 9.9$): thin markets with tight budgets generate the largest efficiency losses. Value dispersion reduces PoA ($|t| = 4.8$), as heterogeneous valuations create clearer value rankings that the allocator can exploit. All observed PoA values remain well within the theoretical worst-case bound of~2 \citep{Aggarwal2019, Gaitonde2023}, consistent with Experiment~4a. Allocative efficiency follows a similar pattern: the budget multiplier dominates ($|t| = 44.7$), with loose budgets yielding near-efficient allocations.

\subsection{Revenue and Budget Dynamics}

The auction type $\times$ budget multiplier interaction is the strongest two-way effect on revenue ($|t| = 9.5$), indicating that the revenue advantage of first-price auctions is concentrated in loose-budget settings where agents can bid aggressively. Under tight budgets, both formats converge as pacing constraints bind. The auction type $\times$ value dispersion interaction ($|t| = 5.6$) reflects a similar pattern: first-price auctions benefit disproportionately from high-$\sigma$ environments, where heterogeneous valuations create wider bid spreads. The budget multiplier $\times$ value dispersion ($|t| = 4.2$) and number of bidders $\times$ value dispersion ($|t| = 3.5$) interactions round out the significant two-way effects. Figure~\ref{fig:e4b_int_rev} displays these interaction patterns.

\begin{figure}[H]
  \centering
  \includegraphics[width=0.7\textwidth]{figures/e4b_int_rev}
  \caption{Experiment~4b: Interaction plot for platform revenue.}
  \label{fig:e4b_int_rev}
\end{figure}

\subsection{Bidding Behaviour}

The bid-to-value ratio is well-explained by the factorial model ($R^2 = 0.89$). The budget multiplier is the dominant factor ($|t| = 54.4$): loose budgets permit higher bids relative to value, while tight budgets force aggressive shading. Auction format ranks second ($|t| = 22.3$), with first-price auctions producing substantially lower bid-to-value ratios than second-price, consistent with strategic bid shading under pay-as-bid rules. The auction type $\times$ budget multiplier interaction ($|t| = 17.3$) indicates that bid shading under first-price rules is amplified when budgets are tight, as pacing and strategic incentives reinforce each other. The auction type $\times$ number of bidders interaction ($|t| = 5.8$) reflects that additional bidders narrow the bid-to-value gap between formats. Compared with the multiplicative pacing of Experiment~4a, the PI controller produces qualitatively similar bid-to-value patterns but with different magnitudes: the additive update rule and cumulative error tracking of the PI controller yield more gradual bid adjustments, resulting in less extreme shading under tight budgets. Table~\ref{tab:exp4b_ranked_vol} and Figure~\ref{fig:e4b_main_vol} report the full effect hierarchy.

\input{tables/exp4b_ranked_vol}

\begin{figure}[H]
  \centering
  \includegraphics[width=0.7\textwidth]{figures/e4b_main_vol}
  \caption{Experiment~4b: Main effects plot for bid-to-value ratio.}
  \label{fig:e4b_main_vol}
\end{figure}

\subsection{Summary}

The comparison between multiplicative dual pacing and PI control isolates the role of the control mechanism from the budget constraint itself. As in Experiment~4a, the budget multiplier is the dominant factor ($|t| = 31.2$), with the number of bidders ranking second ($|t| = 16.1$) and auction format third ($|t| = 11.0$). All PoA values remain well within the theoretical worst-case bound of~2. The PI controller's additive updates and cumulative error tracking produce different bid-to-value ratio levels than multiplicative pacing, but the factor hierarchy is consistent across both approaches.
